\section*{Preface: what this document is, and how to read it}
\addcontentsline{toc}{section}{Preface}

This monograph is not a single theorem with a proof. It is the map of a research programme.
Covering the unit disk with $100$ equal disks is a continuous, non-convex problem whose
combinatorial type changes with the radius, and no known method solves it in one step. What exists
instead is a set of \emph{routes}: constructions that produce upper bounds, relaxations that produce
lower bounds, and a large number of variants that were tried and either closed a sub-case, produced
a diagnostic, or were falsified. The purpose of this document is to present those routes
systematically, so that a reader can see (i) what the current frontier is, (ii) which inequality
carries which part of it, and (iii) exactly where the remaining gap lives.

\subsection*{The one-page state of the art}

\begin{center}
\footnotesize
\begin{tabular}{@{}>{\raggedright\arraybackslash}p{0.30\textwidth}>{\raggedright\arraybackslash}p{0.60\textwidth}@{}}
\toprule
Item & Value / status \\
\midrule
Accepted strict interval &
$0.113235<\rD(100)\le 0.1154469548855829$ \\[1mm]
Accepted lower bound (exact rational) &
$R_\star=22647/200000=0.113235$ \\[1mm]
Accepted upper bound (exact rational) &
$U_\star=1154469548855829/10^{16}$ \\[1mm]
Relative gap $(U_\star-R_\star)/U_\star$ &
approximately $1.92\%$ (display only) \\[1mm]
Lower-bound engine &
six mutually exclusive boundary-count branches
$B31,B32,B33,B34,B35,B{\ge}36$; $B31$--$B35$ inherited from $R=2831/25000$ by radius
monotonicity, $B{\ge}36$ replayed at $R_\star$ \\[1mm]
Weakest branch &
$B\ge 36$, specifically the single value $B=36$ \\[1mm]
Target not reached &
$0.11324$; the frozen capped-pair relaxation has an exact obstruction point at
$x=y=9921033/12500000$ \\[1mm]
Undecided window &
$(0.113235,\,0.11324]$; the rational candidates $0.113236$ and $0.113237$ were started and
stopped without producing a certificate \\[1mm]
Exact optimum, minimal polynomial &
not obtained; \texttt{global\_\allowbreak{}optimality\_\allowbreak{}proved = false} \\[1mm]
Best universal (configuration-free) lower bound &
$2787/25000=0.11148$ (long-chord bundle), with the older
$554570693/5000000000=0.1109141386$ and
$1111720153/10^{10}=0.1111720153$ still valid and dominated \\
\bottomrule
\end{tabular}
\end{center}

\subsection*{How the parts fit together}

\begin{itemize}[leftmargin=1.4em,itemsep=2pt]
\item \textbf{Part I} fixes the problem, the conventions (closed disks, centres unrestricted,
rational decision layer) and the evidence discipline. Read this first: the words
\emph{certificate}, \emph{diagnostic}, \emph{analytic} and \emph{inherited} have precise meanings
here.
\item \textbf{Part II} collects the upper-bound routes: how a finite check replaces a continuous
coverage statement, the triangular-lattice $97{+}3$ configuration with $\alpha=1/\sqrt{67}$, the
replayed external $100$-centre configuration, the cooperative search track, and the local
algebraic (KKT) route.
\item \textbf{Part III} collects the configuration-free analytic lower bounds: the area bound
$1/10$, the integer-edge-count route with its four geometric lemmas, the sharpened envelope, the
long-chord quadratic-loss bound, and the refutation of the perimeter-defect ($\lambda$) law.
\item \textbf{Part IV} is the technical core: the six-branch engine. It explains the same-radius
normalization, the counting bridge, and then each branch $B31$--$B35$ and the tail $B\ge36$,
including the capped-pair relaxation and its rational square-domain subdivision.
\item \textbf{Part V} is the frontier: how the accepted bound moved from $0.11323$ to $0.1132301$
to $0.113235$, why $0.11324$ is blocked, and what a genuinely new inequality would have to do.
\item \textbf{Part VI} is the route catalogue: every route family with its construction, its
status and its evidence path, including the falsified and zero-yield ones.
\item \textbf{Part VII} gives reproduction commands, external dependencies and the boundary of
what was absorbed from third-party material.
\item \textbf{Part VIII} lists the open obligations without softening them.
\item \textbf{The Appendix} is the evidence ledger: assertion $\leftrightarrow$ artifact, plus the
JSON field vocabulary and the exact arithmetic of the $B=36$ obstruction.
\end{itemize}

\subsection*{Evidence grades used throughout}

\begin{center}
\footnotesize
\begin{tabular}{@{}c>{\raggedright\arraybackslash}p{0.50\textwidth}>{\raggedright\arraybackslash}p{0.20\textwidth}@{}}
\toprule
Grade & Meaning & Trust \\
\midrule
A & structured machine artifact: JSON certificate with status, counts, hashes, exit code &
strongest \\
B & independent replay: a different script or stage rechecks the same sub-assertion & strong \\
C & discovery material: floating-point search, candidate ranking, failed experiment &
guides the next step only \\
D & inheritance reasoning, e.g. exclusion at a larger radius used at a smaller radius &
valid only if recorded explicitly by the combiner \\
\bottomrule
\end{tabular}
\end{center}

Three further qualifiers appear constantly and must not be confused with grades:

\begin{itemize}[leftmargin=1.4em,itemsep=2pt]
\item \emph{analytic} --- proved in prose, human-checkable, but not formalized in a proof
assistant. Every ``universal'' statement in this document is analytic plus exact finite arithmetic;
none is a kernel-verified theorem.
\item \emph{\texttt{machine\_\allowbreak{}verified\_\allowbreak{}geometry = false}} --- the finite rational/interval
inequalities were machine-checked; the reduction from geometry to those inequalities was not.
\item \emph{diagnostic} --- the computation deliberately reports a failure, a bottleneck or a
zero-yield result and is \emph{not} a certificate that a branch is closed.
\end{itemize}

\subsection*{Notation}

\begin{center}
\footnotesize
\begin{tabular}{@{}ll@{}}
\toprule
Symbol & Meaning \\
\midrule
$D$ & closed unit disk $\{x\in\R^2:\|x\|\le 1\}$ \\
$B(c,r)$ & closed disk of radius $r$ about $c$ \\
$\rD(n)$ & $\inf\{r: D\subseteq\bigcup_{i=1}^{n}B(c_i,r)\}$ \\
$R$ & a fixed rational \emph{candidate} radius used for a contradiction argument \\
$U$ & a certified upper bound (radius of an explicit covering configuration) \\
$B$ & number of boundary arcs (equivalently boundary vertices) served by the covering \\
$M$ & number of interior faces of the normalized common triangulation \\
$N$, $m$ & number of distinct centres, resp.\ number of effectively used disks \\
$k$ & number of distinct boundary-serving disks \\
$e_i$ & $m_i+a_i$, the nominal side count of the clipped cell $Q_i$ \\
$\Gamma$ & total marked interior angle sum used in the angular ledgers \\
$h$ & total degree excess of strictly deep vertices \\
$C$ & number of chords of a core boundary \\
$\lambda,\mu,L,K$ & penalty coefficient, slope, intercept, marked constant of the tail relaxation \\
$\Phi$ & local pair loss of the capped-pair relaxation \\
\bottomrule
\end{tabular}
\end{center}

\subsection*{A warning about reading the numbers}

Every rational number in this document is quoted as it appears in a certificate or a report. Where
a decimal is shown for orientation it is marked as display. Where two stages give slightly
different auxiliary constants (for example the normalized centre-norm-squared bound at different
radii), both are quoted with their own source rather than silently reconciled. Where a run was
started and stopped, it is reported as stopped --- not as a failure and not as a proof.
